\chapter{Net worth and aggregation}

\section{What was missing}

An early version of the application tracked investments and called their value
"net worth". That was wrong in an obvious way: it ignored the money sitting in
bank accounts, which for most people is a large part of what they have.

Adding accounts was easy. Deciding how they participate in the rest of the
application was the real work.

\section{Accounts}

An account is a name, a kind (account, savings or cash), a balance and a
timestamp. There is no bank integration, and that is a decision rather than a
limitation: connecting to a bank means credentials, aggregation providers and a
recurring cost, for data the user can update in ten seconds a month.

Every write to an account ends the same way:

\begin{lstlisting}[style=py]
db.commit()
snapshot_today(db)
\end{lstlisting}

So the evolution chart follows real changes instead of only moving when a
statement is imported.

\section{The aggregation module}

\file{app/services/aggregation.py} is the only place that computes totals. The
web pages, the PDF reports and the Telegram bot all read from it, which is what
guarantees that the same number appears everywhere.

Its main function builds the whole dashboard payload in one pass:

\begin{lstlisting}[style=py]
def build_dashboard(db, year=None, month=None) -> dict:
    invs = latest_investments(db)
    inv_value = sum(_f(i.current_value) for i in invs)
    inv_invested = sum(_f(i.invested) for i in invs)
    inv_profit = sum(_f(i.profit) for i in invs)
    return_pct = (inv_profit / inv_invested * 100) if inv_invested else 0.0

    banks = db.scalars(select(BankAccount).order_by(BankAccount.balance.desc())).all()
    bank_value = sum(_f(b.balance) for b in banks)
    net_worth = inv_value + bank_value
    ...
    return {"year": year, "month": month, "kpis": {...}, "charts": {...}}
\end{lstlisting}

Note the guard on the division. An empty portfolio has zero invested, and a
return of zero is the honest answer; raising or returning infinity would be
worse.

\section{Month ranges without special cases}

Several features ask for a range of months: a report from October 2025 to March
2026, a cashflow chart, a net figure for any period. Written naively, those
queries fill up with conditions about years and boundaries.

The module converts a year and a month into a single integer, the month ordinal:

\begin{lstlisting}[style=py]
def _month_ordinals(y1, m1, y2, m2) -> tuple[int, int]:
    start, end = y1 * 12 + m1, y2 * 12 + m2
    return (end, start) if start > end else (start, end)

_ORDINAL = Transaction.year * 12 + Transaction.month
\end{lstlisting}

A range then becomes a plain comparison, \code{\_ORDINAL >= start} and
\code{\_ORDINAL <= end}, with no year logic at all. The same expression is built
once by SQLAlchemy and reused by every range query.

Reversed inputs are normalised in the same helper, so a user who types the end
month first gets the range they meant instead of an empty result.

Turning an ordinal back into a year and a month is one line:

\begin{lstlisting}[style=py]
for o in range(start, end + 1):
    y, m0 = divmod(o - 1, 12)
    months.append((y, m0 + 1))
\end{lstlisting}

\begin{keypoint}[The general lesson]
When a range keeps producing boundary bugs, it is usually the representation
that is wrong, not the conditions. Mapping the values onto a single ordered
number makes the boundaries disappear.
\end{keypoint}

\section{The dashboard}

The dashboard answers two questions on one screen: how much does the user have,
and where did this month go.

Four indicators sit at the top: net worth (with the split between investments
and banks, plus profit and return), income of the month, expenses of the month,
and savings of the month with the invested part noted underneath.

Four charts follow:

\begin{itemize}
  \item income against expenses per month of the year, as bars;
  \item expenses by category for the selected month, as a doughnut;
  \item net worth over time, as a line, from the snapshot table;
  \item the split of net worth by broker and account, as a doughnut.
\end{itemize}

The net worth series has a fallback. When there are no snapshots yet, because
the application was just installed, it groups investments by valuation date
instead, so the chart is not empty on the first day.

\begin{lstlisting}[style=py]
if snaps:
    evo_labels = [d.strftime("%d/%m/%y") for d, _ in snaps]
    evo_data = [round(_f(s), 2) for _, s in snaps]
else:
    evo = db.execute(select(Investment.valued_on, func.sum(Investment.current_value))
                     .group_by(Investment.valued_on)
                     .order_by(Investment.valued_on)).all()
\end{lstlisting}

\section{An empty database is a valid state}

Every aggregation function is written so that a fresh installation renders
correctly. Sums return zero instead of \code{None}, thanks to a small helper:

\begin{lstlisting}[style=py]
def _f(x) -> float:
    return float(x or 0)
\end{lstlisting}

The list of years falls back to the current year when there are no transactions,
so the year selector is never empty. And the payload carries a \code{has\_data}
flag, which the template uses to show a short message pointing at the two places
where a new user should start.

It is a small thing, and it is the difference between an application that looks
broken on day one and one that looks ready.
